\begin{proof}[Proof of \cref{lem:p142-monotone-parameter-supremum}]
For $N\geq1$ put $a=c-N$, and partition $[a,c)$ into the $N^2$
half-open intervals
\[
 [\tau_{N,j},\tau_{N,j}+1/N),\qquad
 \tau_{N,j}=c-N+j/N,\quad 0\leq j<N^2.
\]
Only left endpoints are sampled; the family is never evaluated at $c$.
Set
\[
 F_N:=\max\{c-N,F\},\qquad
 G_N:=\max\left\{c-N,
      \max_{0\leq j<N^2}
       (\tau_{N,j}-a_{\tau_{N,j}})\right\}.
\]
Every sampled term occurs in the supremum defining $F$, so $G_N\leq F_N$.
If $\tau<c-N$, nonnegativity gives $\tau-a_\tau\leq c-N$.
Otherwise there is $j$ with
$\tau_{N,j}\leq\tau<\tau_{N,j}+1/N$. Monotonicity in the parameter
gives, in extended-real order,
\[
 \tau-a_\tau
 \leq \tau-a_{\tau_{N,j}}
 \leq \tau_{N,j}-a_{\tau_{N,j}}+1/N
 \leq G_N+1/N.
\]
This also holds when one or both of the ideal-type quantities are infinite:
only a finite real number minus $+\infty$ occurs. Taking the supremum and
including the floor proves
\begin{equation}
             G_N\leq F_N\leq G_N+1/N.
 \label{eq:p142-grid-sandwich}
\end{equation}

Each $G_N$ is a finite-valued upper-semicontinuous function, since
$-a_{\tau_{N,j}}$ is upper semicontinuous. The sandwich gives the exact
identity
\begin{equation}
                  F=\inf_{N\geq1}(G_N+1/N).
 \label{eq:p142-usc-infimum}
\end{equation}
Indeed, $F\leq F_N\leq G_N+1/N$, and
$G_N+1/N\leq F_N+1/N$ tends pointwise to $F$ as $N\to\infty$.
If $F(y)=-\infty$, then $F_N(y)=c-N$, so the same conclusion holds
at that point. An arbitrary infimum of upper-semicontinuous functions is
upper semicontinuous: its strict sublevel set is the union of the strict
sublevel sets of its terms. This proves the first assertion on an
arbitrary topological space.

For the dense-parameter assertion, fix $y$ and $\tau<c$ with
$a_\tau(y)<+\infty$. Given $\varepsilon>0$, choose
$s\in D\cap(\tau-\varepsilon,\tau)$. Then
\[
 s-a_s(y)\geq s-a_\tau(y)
              >\tau-a_\tau(y)-\varepsilon.
\]
Thus the supremum over $D$ is at least every finite term in the full
supremum, after letting $\varepsilon$ tend to zero. Terms equal to
$-\infty$ impose no further condition. If all terms are $-\infty$,
both suprema are $-\infty$. The opposite inequality follows from
$D\subset(-\infty,c)$, proving
\eqref{eq:p142-dense-parameter-formula}.
\end{proof}

\begin{proof}[Proof of \cref{thm:p142-raw-transform-usc}]
Multiplier ideal sheaves of quasi-psh functions on a complex manifold
are coherent. Their subadditivity \cite[Theorem~2.6(ii)]{DEL00} gives
\[
 \mathcal I((m+\ell)\Gamma_\tau)
 \subseteq\mathcal I(m\Gamma_\tau)
               \mathcal I(\ell\Gamma_\tau).
\]
Ideal values reverse inclusion and add under multiplication. Hence
$b_m=v(\mathcal I(m\Gamma_\tau))$ is a nonnegative extended-real
superadditive sequence. If some $b_q=+\infty$, superadditivity implies
$b_m=+\infty$ for every $m\geq q$: use $m=q+r$ for $r\geq1$.
Otherwise, for a fixed $q$, writing $m=kq+r$, with $0\leq r<q$, gives
$b_m\geq kb_q$ (put $b_0=0$). Therefore
\[
 \liminf_{m\to\infty}\frac{b_m}{m}\geq\frac{b_q}{q}.
\]
Taking the supremum over $q$ and comparing with the immediate upper
bound proves the limit assertion, including an infinite supremum. This
is the multiplier-ideal argument recorded in
\cref{prop:pseff-ideal-valuative-basic}.

By the topology of the tropical spectrum, the function
$v\mapsto v(\mathfrak a)$ is continuous into $[0,+\infty]$ with its
order topology for every coherent ideal $\mathfrak a$; see
\cite[Definition~1.3.1]{MP24} and \eqref{eq:ts-evaluation}.
It follows that
\[
 v\longmapsto\nu_v^{\mathcal I}(\Gamma_\tau)
   =\sup_{m\geq1}\frac1m
          v(\mathcal I(m\Gamma_\tau))
\]
is lower semicontinuous. If $s\leq t<c$, then
$\Gamma_s\geq\Gamma_t$. Local integrability immediately implies
\[
 \mathcal I(m\Gamma_t)\subseteq\mathcal I(m\Gamma_s),
 \qquad
 \nu_v^{\mathcal I}(\Gamma_s)
       \leq\nu_v^{\mathcal I}(\Gamma_t).
\]
For the first inclusion, multiplication of the weight by the bounded
factor associated to an additive local constant would also suffice;
the given pointwise order already yields it directly. No assertion
about continuity of multiplier ideals in the parameter is used.
The hypotheses of \cref{lem:p142-monotone-parameter-supremum} are now
verified, and that lemma proves upper semicontinuity and the dense-set
formula at every semivaluation.

We prove the stronger explicit approximation assertion separately.
For fixed $N,M$, all ideal evaluations in
\eqref{eq:p142-upper-approximants} are continuous extended-real
functions. Their displayed finite maxima and positive scalings are
continuous in the order topology; subtracting from the finite
$\tau_{N,j}$ gives a continuous function into $[-\infty,\tau_{N,j}]$.
The outer maximum with $c-N$ then yields an ordinary real-valued
continuous function. In particular,
\[
        c-N+1/N\leq B_{N,M}\leq c+1/N.
\]
This establishes boundedness without invoking compactness of $X^\beth$.

As $M$ increases, $A_{\tau,M}\nearrow\nu^{\mathcal I}(\Gamma_\tau)$.
For the finite number of sampled parameters, taking the maximum commutes
with these decreasing limits after changing the sign. To see this even
at $-\infty$, if every limiting term is below a chosen real threshold,
then every term is eventually below it; a maximum of finitely many
indices works for all terms. The fixed finite floor is retained throughout.
Consequently, with $G_N$ as in the preceding proof,
\[
           B_{N,M}\searrow G_N+1/N\quad(M\to\infty).
\]
The sandwich \eqref{eq:p142-grid-sandwich} gives
\begin{equation}
 F_\Gamma\leq B_{N,M},\qquad
 \inf_{N,M\geq1}B_{N,M}
      =\inf_{N\geq1}(G_N+1/N)=F_\Gamma.
 \label{eq:p142-double-infimum}
\end{equation}
Every $f_k$ is a finite minimum of bounded continuous real-valued
functions. The squares $\{1,\ldots,k\}^2$ increase and exhaust
$\mathbf Z_{>0}^2$, so their minima decrease to the double infimum
in \eqref{eq:p142-double-infimum}. This is
\eqref{eq:p142-exact-sequential-limit}.

The construction is explicit. It does not extract a sequence from an
arbitrary net on a nonmetrizable space. Nor does a minimum of these
continuous functions acquire a positivity property merely because it
is formed from ideal values. No positive approximation is asserted.
\end{proof}

\begin{proof}[Proof of \cref{cor:p142-pseudoeffective-raw-usc}]
Fix one K\"ahler perturbation component
$\Gamma=(\Gamma^{\theta+\omega}_\tau)_{\tau<c}$, where
$c=\boldsymbol\Gamma_{\max}$. The DXZ test-curve definition supplies
a decreasing family of quasi-psh functions with finite terminal
parameter, and its interior levels are not identically $-\infty$;
see \cite[Section~2.2, Definition~3.3]{DXZ23}.
\Cref{thm:p142-raw-transform-usc} therefore applies to this component.

The identification with the raw inverse-system transform uses only the
already established transition identity, not an MP comparison theorem.
Under a transition the interior level is replaced by its
$\mathcal I$-model envelope in the larger class. This preserves every
integer multiplier ideal by \cite[Proposition~2.18(ii)]{DX22}; the
standing hypotheses of its Section~2.4 allow an arbitrary compact
K\"ahler manifold and a pseudoeffective background class. A common
dominating component compares any
two choices, and their terminal parameters coincide. Evaluating those
same ideals at any semivaluation proves that both the raw transform and
the explicit expressions \eqref{eq:p142-finite-ideal-values}--
\eqref{eq:p142-decreasing-sequence} are unchanged. These are precisely
the ideal equalities in
\cref{thm:pseff-perturbation-independent-transform}; the source
transition is \cite[formula~(28), Definition~3.3]{DXZ23}.

All arguments take place on a perturbed component. No potential of
positive mass in the unperturbed class is chosen, so zero-volume
pseudoeffective classes are included. The proof does not construct
relatively nef cohomological models, a Monge--Amp\`ere operator, or an
energy. Those are distinct questions.
\end{proof}
